% Generated by src/equivalent_rights_proposals.jl -- do not edit by hand.
\begin{table}[htbp]
\centering
\small
\caption{Club-price rights ratios $\rho_i$ equivalent to the Wolfram et al. and Banerjee et al. proposals}
\renewcommand{\arraystretch}{1.15}
\begin{tabular}{lcccccc}
  \toprule
  & \multicolumn{3}{c}{\textbf{Wolfram et al.}} & \multicolumn{3}{c}{\textbf{Banerjee et al.}} \\
  \cmidrule(lr){2-4} \cmidrule(lr){5-7}
  \textbf{Country} & $p_i$ & $\rho^{A}$ & $\rho^{B}$ & $p_i$ & $\rho^{A}$ & $\rho^{B}$ \\
  \midrule
  United States & 0 & -- & -- & 75 & 2.16 & 2.82 \\
  Russia & 0 & -- & -- & 75 & 1.36 & 1.89 \\
  European Union & 75 & 0.26 & 0.36 & 75 & 0.26 & 0.36 \\
  Turkey & 0 & -- & -- & 50 & 1.54 & 1.62 \\
  Indonesia & 50 & 0.89 & 0.92 & 50 & 1.09 & 1.04 \\
  China & 50 & 0.84 & 0.98 & 50 & 0.88 & 0.95 \\
  Brazil & 50 & 0.38 & 0.37 & 50 & 0.5 & 0.41 \\
  India & 25 & 1.3 & 1.12 & 30 & 1.18 & 1.03 \\
  Nigeria & 0 & -- & -- & 30 & 0.24 & 0.11 \\
  DR Congo & 0 & -- & -- & 10 & 0.06 & 0.01 \\
  \midrule
  \textbf{Gain in rights, variant 1 (\%)} &  & \rose{9.9\%} & \rose{9.5\%} &  & \rose{9.1\%} & \rose{8.5\%} \\
  \textbf{World temp.~2100, change (\textdegree{}C)} &  & \rose{$-$0.020} & \rose{$-$0.019} &  & \rose{$-$0.037} & \rose{$-$0.034} \\
  \textbf{World welfare gain, EDE (variant 2)} &  & 0.028\% & $-$0.006\% &  & 0.102\% & $-$0.009\% \\
  \textbf{World consumption gain (variant 2)} &  & 0.022\% & 0.018\% &  & 0.057\% & 0.048\% \\
  \textbf{World welfare gain, EDE (variant 3)} &  & 0.028\% & 0.011\% &  & 0.102\% & 0.041\% \\
  \textbf{World consumption gain (variant 3)} &  & 0.022\% & 0.020\% &  & 0.057\% & 0.047\% \\
  \textbf{World welfare gain, EDE (variant 4)} &  & 0.081\% & 0.048\% &  & 0.419\% & 0.298\% \\
  \textbf{World consumption gain (variant 4)} &  & 0.024\% & 0.021\% &  & 0.068\% & 0.059\% \\
  \bottomrule
\end{tabular}
\label{tab:equiv_rights}
\\[4pt]
{\footnotesize Note: $p_i$ is the carbon price the proposal asks of country $i$ in 2030; the comparison regime instead prices every member of the proposal's own club at one common price (\$49.2/t for Wolfram et al. and \$55.8/t for Banerjee et al. in 2030) and leaves non-members unpriced, as the proposal does. $\rho_i$ is the country's allocation as a multiple of an equal-per-capita share of the club's emissions under the proposal, taken at face value in variant 1; $\rho^{A}$ solves each country on its own, $\rho^{B}$ all of them jointly. The European Union figure aggregates its members' rights, and a country the proposal does not price is outside the club in both regimes, shown at $p_i=0$ with no equivalent allocation (`--'). Variant 1 lets the equivalent allocation set the cap; variants 2 to 4 return the resulting surplus as rights so that club emissions match the proposal's, sharing it by uniform scaling (2), after a floor on option A's allocation (3), and in proportion to population times marginal utility (4). The temperature row is the change from the proposal's own 2100 warming, which is 2.16\textdegree{}C for Wolfram et al. and 1.80\textdegree{}C for Banerjee et al.. Welfare is reported on the inequality-averse global EDE ($\eta=1.5$) and on world mean consumption per capita, which weights every person's consumption equally.}
\end{table}
